\documentclass{article}
% notation.tex --- shared preamble for Warp paper-mods.
%
% Every mod*.tex file begins with:
%
%     % notation.tex --- shared preamble for Warp paper-mods.
%
% Every mod*.tex file begins with:
%
%     % notation.tex --- shared preamble for Warp paper-mods.
%
% Every mod*.tex file begins with:
%
%     \input{notation.tex}
%
% Conventions
% -----------
% * Framing: IOP in the sense of Ben-Sasson--Chiesa--Spooner (BCS, TCC 2016),
%   NOT Algebraic Holographic Proofs. Oracles are functions $f:[n] \to \F$,
%   queried by index; multilinear-extension semantics layer on top for point
%   queries.
% * Every modification file cites the Rust module path it pairs with.

\usepackage{amsmath,amssymb,amsthm}
\usepackage{mathtools}
\usepackage{hyperref}

% ---- Fields, indexing, domains -------------------------------------------
\newcommand{\F}{\mathbb{F}}
\newcommand{\idx}[1]{\ensuremath{[#1]}}               % index set [n]
\newcommand{\bool}[1]{\ensuremath{\{0,1\}^{#1}}}      % boolean hypercube
\newcommand{\size}[1]{\ensuremath{\left\lvert#1\right\rvert}}

% ---- Oracles (Modification 1) --------------------------------------------
% An Oracle carries three views of the same object:
%   (i)   an evaluation table  f : [n] -> F        (BCS-native)
%   (ii)  a multilinear extension  \hat f : F^{log n} -> F
%   (iii) a commitment handle (Merkle root under BCS compilation)
% See mod1_oracle.tex.
\newcommand{\Oracle}[1]{\ensuremath{\mathsf{O}[#1]}}
\newcommand{\MLE}[1]{\ensuremath{\widehat{#1}}}
\newcommand{\Commit}[1]{\ensuremath{\mathsf{com}(#1)}}
\newcommand{\query}{\ensuremath{\mathsf{query}}}      % generic query operator

% ---- Accumulator state (Modification 3, reserved) -------------------------
\newcommand{\AccX}{\ensuremath{\mathsf{acc}_x}}       % accumulator instance
\newcommand{\AccW}{\ensuremath{\mathsf{acc}_w}}       % accumulator witness
\newcommand{\AccState}{\ensuremath{(\AccX, \AccW)}}

% ---- Common structural IORs ----------------------------------------------
\newcommand{\IOR}[1]{\textsc{#1}}
\newcommand{\Zerocheck}{\IOR{Zerocheck}}
\newcommand{\Sumcheck}{\IOR{Sumcheck}}
\newcommand{\Pesat}{\IOR{Pesat}}
\newcommand{\TwinConstraint}{\IOR{TwinConstraint}}
\newcommand{\Proximity}{\IOR{Proximity}}
\newcommand{\Batching}{\IOR{Batching}}
\newcommand{\OOD}{\IOR{OOD}}

% ---- Transcript notation --------------------------------------------------
\newcommand{\absorb}{\ensuremath{\lhd}}               % absorb into transcript
\newcommand{\squeeze}{\ensuremath{\rhd}}              % derive from transcript

% ---- Code / module cross-reference ---------------------------------------
% Use \codemod{src/protocol/oracle.rs} in .tex files to point at a Rust
% module. Renders as inline typewriter with a URL when compiled with hyperref.
\newcommand{\codemod}[1]{\href{../../#1}{\texttt{#1}}}

%
% Conventions
% -----------
% * Framing: IOP in the sense of Ben-Sasson--Chiesa--Spooner (BCS, TCC 2016),
%   NOT Algebraic Holographic Proofs. Oracles are functions $f:[n] \to \F$,
%   queried by index; multilinear-extension semantics layer on top for point
%   queries.
% * Every modification file cites the Rust module path it pairs with.

\usepackage{amsmath,amssymb,amsthm}
\usepackage{mathtools}
\usepackage{hyperref}

% ---- Fields, indexing, domains -------------------------------------------
\newcommand{\F}{\mathbb{F}}
\newcommand{\idx}[1]{\ensuremath{[#1]}}               % index set [n]
\newcommand{\bool}[1]{\ensuremath{\{0,1\}^{#1}}}      % boolean hypercube
\newcommand{\size}[1]{\ensuremath{\left\lvert#1\right\rvert}}

% ---- Oracles (Modification 1) --------------------------------------------
% An Oracle carries three views of the same object:
%   (i)   an evaluation table  f : [n] -> F        (BCS-native)
%   (ii)  a multilinear extension  \hat f : F^{log n} -> F
%   (iii) a commitment handle (Merkle root under BCS compilation)
% See mod1_oracle.tex.
\newcommand{\Oracle}[1]{\ensuremath{\mathsf{O}[#1]}}
\newcommand{\MLE}[1]{\ensuremath{\widehat{#1}}}
\newcommand{\Commit}[1]{\ensuremath{\mathsf{com}(#1)}}
\newcommand{\query}{\ensuremath{\mathsf{query}}}      % generic query operator

% ---- Accumulator state (Modification 3, reserved) -------------------------
\newcommand{\AccX}{\ensuremath{\mathsf{acc}_x}}       % accumulator instance
\newcommand{\AccW}{\ensuremath{\mathsf{acc}_w}}       % accumulator witness
\newcommand{\AccState}{\ensuremath{(\AccX, \AccW)}}

% ---- Common structural IORs ----------------------------------------------
\newcommand{\IOR}[1]{\textsc{#1}}
\newcommand{\Zerocheck}{\IOR{Zerocheck}}
\newcommand{\Sumcheck}{\IOR{Sumcheck}}
\newcommand{\Pesat}{\IOR{Pesat}}
\newcommand{\TwinConstraint}{\IOR{TwinConstraint}}
\newcommand{\Proximity}{\IOR{Proximity}}
\newcommand{\Batching}{\IOR{Batching}}
\newcommand{\OOD}{\IOR{OOD}}

% ---- Transcript notation --------------------------------------------------
\newcommand{\absorb}{\ensuremath{\lhd}}               % absorb into transcript
\newcommand{\squeeze}{\ensuremath{\rhd}}              % derive from transcript

% ---- Code / module cross-reference ---------------------------------------
% Use \codemod{src/protocol/oracle.rs} in .tex files to point at a Rust
% module. Renders as inline typewriter with a URL when compiled with hyperref.
\newcommand{\codemod}[1]{\href{../../#1}{\texttt{#1}}}

%
% Conventions
% -----------
% * Framing: IOP in the sense of Ben-Sasson--Chiesa--Spooner (BCS, TCC 2016),
%   NOT Algebraic Holographic Proofs. Oracles are functions $f:[n] \to \F$,
%   queried by index; multilinear-extension semantics layer on top for point
%   queries.
% * Every modification file cites the Rust module path it pairs with.

\usepackage{amsmath,amssymb,amsthm}
\usepackage{mathtools}
\usepackage{hyperref}

% ---- Fields, indexing, domains -------------------------------------------
\newcommand{\F}{\mathbb{F}}
\newcommand{\idx}[1]{\ensuremath{[#1]}}               % index set [n]
\newcommand{\bool}[1]{\ensuremath{\{0,1\}^{#1}}}      % boolean hypercube
\newcommand{\size}[1]{\ensuremath{\left\lvert#1\right\rvert}}

% ---- Oracles (Modification 1) --------------------------------------------
% An Oracle carries three views of the same object:
%   (i)   an evaluation table  f : [n] -> F        (BCS-native)
%   (ii)  a multilinear extension  \hat f : F^{log n} -> F
%   (iii) a commitment handle (Merkle root under BCS compilation)
% See mod1_oracle.tex.
\newcommand{\Oracle}[1]{\ensuremath{\mathsf{O}[#1]}}
\newcommand{\MLE}[1]{\ensuremath{\widehat{#1}}}
\newcommand{\Commit}[1]{\ensuremath{\mathsf{com}(#1)}}
\newcommand{\query}{\ensuremath{\mathsf{query}}}      % generic query operator

% ---- Accumulator state (Modification 3, reserved) -------------------------
\newcommand{\AccX}{\ensuremath{\mathsf{acc}_x}}       % accumulator instance
\newcommand{\AccW}{\ensuremath{\mathsf{acc}_w}}       % accumulator witness
\newcommand{\AccState}{\ensuremath{(\AccX, \AccW)}}

% ---- Common structural IORs ----------------------------------------------
\newcommand{\IOR}[1]{\textsc{#1}}
\newcommand{\Zerocheck}{\IOR{Zerocheck}}
\newcommand{\Sumcheck}{\IOR{Sumcheck}}
\newcommand{\Pesat}{\IOR{Pesat}}
\newcommand{\TwinConstraint}{\IOR{TwinConstraint}}
\newcommand{\Proximity}{\IOR{Proximity}}
\newcommand{\Batching}{\IOR{Batching}}
\newcommand{\OOD}{\IOR{OOD}}

% ---- Transcript notation --------------------------------------------------
\newcommand{\absorb}{\ensuremath{\lhd}}               % absorb into transcript
\newcommand{\squeeze}{\ensuremath{\rhd}}              % derive from transcript

% ---- Code / module cross-reference ---------------------------------------
% Use \codemod{src/protocol/oracle.rs} in .tex files to point at a Rust
% module. Renders as inline typewriter with a URL when compiled with hyperref.
\newcommand{\codemod}[1]{\href{../../#1}{\texttt{#1}}}


\title{Modification 4: \\ Soundness-driven parameter selection}
\author{Warp paper-mods}
\date{}

\begin{document}
\maketitle

\paragraph{Paired Rust module.} \codemod{src/params/}, with a CLI front-end
at \codemod{src/bin/warp-params.rs}.

\section{Motivation}

The paper states the soundness bound of Warp symbolically; the implementation
has so far hard-coded parameter values (for instance, \(s = 8\), \(t = 7\)
in the regression tests). There is no record of which security target those
choices correspond to, nor is there a way to ask the code for the minimum
parameter tuple that hits a chosen target \(\lambda\).

Modification 4 adds, on both sides, an \emph{inverse} to the soundness
analysis. Given
\(\lambda \in \mathbb{N}\), field size \(\size{\F}\), Reed--Solomon code rate
\(\rho = k/n\), and a choice of \emph{list-decoding regime}, produce the
smallest parameter tuple \((s,\ t)\) that provably (or conjecturally) gives
\(\lambda\) bits of soundness. The workload parameters \((l,\ l_1)\) remain
caller-chosen; Modification 4 does not select them.

\section{Regimes}

\begin{description}
  \item[\IOR{Provable}.] Each proximity query rejects a word at distance
        \(e\) from the nearest codeword with probability at least \(e\).
        Taking the Johnson bound \(e = 1 - \sqrt{\rho}\) as the largest
        value supported by a proof for Reed--Solomon codes, the soundness
        error of \(t\) independent proximity queries is at most
        \(\bigl(1 - e\bigr)^{t} = \rho^{t/2}\). Therefore
        \[
            t \;\geq\; \frac{2\lambda}{\log_2(1/\rho)}.
        \]
  \item[\IOR{Conjectured}.] The list-decodability conjecture for
        Reed--Solomon codes extends the argument up to radius
        \(e = 1 - \rho\). The same calculation gives
        \[
            t \;\geq\; \frac{\lambda}{\log_2(1/\rho)}.
        \]
\end{description}

For \(s\) (OOD samples) we use the conservative constant \(s = 8\),
which dominates the round-complexity of batching without materially
affecting soundness at the scales we run. A tighter bound is deferred to
a follow-up; the current code documents this choice explicitly.

\section{Field-size admissibility}

Several other soundness terms are each at most
\(\mathrm{polylog}(\cdot) / \size{\F}\) (sumcheck rounds, zero-check
randomness absorbs, transcript-squeezed challenges). They are negligible
provided \(\log_2 \size{\F} \geq \lambda + \epsilon\) for a small constant
\(\epsilon\); the code rejects any choice that violates this admissibility.

\section{References used, and scope limits}

The formulas above are the standard proximity-gap / list-decoding bounds
from the Reed--Solomon STARK lineage, in the precise formulation used in
STIR~\cite{STIR} and WHIR~\cite{WHIR}. They are \emph{not} a full
re-derivation of the Warp paper's soundness proof; they capture the
dominant term (proximity queries) and ignore polylogarithmic noise
controlled by the field-size check. The Rust module
\codemod{src/params/} uses the same formulas and marks every constant
it uses, so refining against the paper's exact proof reduces to editing
one table.

\paragraph{Deferred.} Calibration of \(s\) against the batching-sumcheck
soundness, derivation for non-Reed--Solomon codes, and a reference
table of attested parameter tuples with matching proofs.

\begin{thebibliography}{9}
\bibitem{STIR}
  Arnon, Chiesa, Fenzi, Yogev.
  \emph{STIR: Reed--Solomon Proximity Testing with Fewer Queries.}
  CRYPTO 2024.
\bibitem{WHIR}
  Arnon, Chiesa, Fenzi, Yogev.
  \emph{WHIR: Reed--Solomon Proximity Testing with Super-Fast Verification.}
  2024.
\end{thebibliography}

\end{document}
